\documentclass[11pt,a4paper]{article}

% ---- Packages ----
\usepackage[utf8]{inputenc}
\usepackage[T1]{fontenc}
\usepackage{amsmath,amssymb}
\usepackage{graphicx}
\usepackage{booktabs}
\usepackage{hyperref}
\usepackage[margin=1in]{geometry}
\usepackage{natbib}
\usepackage{authblk}
\usepackage{setspace}
\usepackage{xcolor}
\usepackage{microtype}

\onehalfspacing

% ---- Title ----
\title{Noncanonical Inflammasome Assembly Requires Caspase-11 Catalytic Activity and Intra-Molecular Autoprocessing}

\author[1]{Jane A.\ Chen}
\author[1]{Michael R.\ Torres}
\author[2]{Sarah L.\ Okonkwo}
\author[1,2]{David K.\ Brennan\thanks{Corresponding author: d.brennan@example.edu}}

\affil[1]{Department of Immunology, Institute of Molecular Medicine, University of Cambridge, Cambridge, UK}
\affil[2]{Centre for Innate Immunity and Inflammation, University of Melbourne, Melbourne, Australia}

\date{}

\begin{document}
\maketitle

% ==============================================================================
\begin{abstract}
% ==============================================================================
Caspase-11 is the cytosolic sensor for bacterial lipopolysaccharide (LPS) and the initiator of the noncanonical inflammasome pathway, which culminates in gasdermin~D (GSDMD) cleavage and pyroptotic cell death. Although caspase-11 is known to oligomerize into a high-molecular-weight complex upon LPS binding, whether it assembles into visible supramolecular organizing centers (SMOCs) in intact cells and whether its enzymatic activity is required for this assembly remain unresolved. Here we report that LPS transfection induces caspase-11--GFP speck formation in murine macrophages, providing direct visual evidence for a noncanonical inflammasome SMOC. Unexpectedly, both a catalytic-dead mutant (C254A) and an autoprocessing-deficient mutant (D285A) fail to form specks and remain monomeric on size-exclusion chromatography, even in the presence of cytosolic LPS. Consistently, both mutants are deficient in GSDMD cleavage, LDH release, and propidium iodide uptake. These findings revise the prevailing model of inflammasome assembly: rather than being passively recruited to a pre-formed oligomeric scaffold, caspase-11 enzymatic activity and intra-molecular autoprocessing actively drive SMOC nucleation and are therefore prerequisites, not consequences, of noncanonical inflammasome complex formation.

\vspace{0.5em}
\noindent\textbf{Keywords:} caspase-11, noncanonical inflammasome, pyroptosis, gasdermin~D, SMOC, LPS, autoprocessing
\end{abstract}

% ==============================================================================
\section{Introduction}
\label{sec:introduction}
% ==============================================================================

Inflammasomes are cytoplasmic multiprotein complexes that serve as sentinels of the innate immune system, detecting pathogen-associated molecular patterns and danger signals to initiate inflammatory responses~\citep{martinon2002, schroder2010}. Upon activation, these complexes recruit and activate inflammatory caspases through proximity-induced dimerization and autoprocessing, ultimately leading to the maturation of pro-inflammatory cytokines and, in many cases, a lytic form of cell death known as pyroptosis~\citep{broz2016, latz2013}.

The canonical inflammasome pathway, exemplified by the NLRP3 inflammasome, is relatively well characterized: sensor proteins detect specific stimuli, recruit the adaptor ASC, and assemble into large oligomeric structures termed supramolecular organizing centers (SMOCs), visualized microscopically as perinuclear specks~\citep{broz2016, man2017}. Within these SMOCs, caspase-1 undergoes proximity-induced activation and cleaves downstream substrates including pro-IL-1$\beta$, pro-IL-18, and gasdermin~D (GSDMD)~\citep{shi2015, he2015}.

In contrast, the noncanonical inflammasome pathway is initiated by a fundamentally different mechanism. Caspase-11 in mice (caspase-4 and caspase-5 in humans) functions as both the sensor and the effector protease~\citep{kayagaki2011, shi2014, hagar2013}. Caspase-11 directly binds intracellular lipopolysaccharide (LPS) from Gram-negative bacteria via its CARD domain, bypassing the need for a canonical sensor protein or ASC adaptor~\citep{shi2014}. LPS binding triggers caspase-11 oligomerization, leading to GSDMD cleavage, pore formation in the plasma membrane, and pyroptosis~\citep{kayagaki2015, shi2015, liu2016, ding2016}.

While biochemical studies have demonstrated that LPS induces the formation of a high-molecular-weight caspase-11 complex~\citep{shi2014, ross2022}, several fundamental questions remain unanswered. First, does caspase-11 form discrete, microscopically visible specks analogous to ASC specks in canonical inflammasomes? Second, is the catalytic activity of caspase-11 required merely for downstream substrate cleavage, or does it play a role in the assembly of the oligomeric complex itself? The prevailing model posits a sequential pathway in which LPS binding induces passive oligomerization, followed by proximity-induced activation of catalytic activity~\citep{broz2016}. Under this model, catalytic activity is a consequence, not a prerequisite, of complex formation.

Here we challenge this prevailing model through a combination of live-cell imaging and biochemical characterization. Using fluorescently tagged caspase-11 constructs expressed in murine macrophages, we provide the first direct visualization of LPS-induced caspase-11 specks. Critically, we demonstrate that a catalytic-dead mutant (C254A) and an autoprocessing-deficient mutant (D285A) both fail to form specks and remain monomeric, even in the presence of cytosolic LPS. These findings establish that caspase-11 catalytic activity and intra-molecular autoprocessing are active drivers of noncanonical inflammasome SMOC assembly, fundamentally revising our understanding of how this innate immune complex is nucleated.

% ==============================================================================
\section{Related Work}
\label{sec:related}
% ==============================================================================

\subsection{Canonical inflammasome assembly}

The concept of the inflammasome was introduced by Martinon et al., who described a caspase-1-activating complex assembled around NLRP1~\citep{martinon2002}. Subsequent work identified multiple canonical inflammasome sensors---including NLRP3, NLRC4, and AIM2---each recognizing distinct stimuli~\citep{schroder2010, latz2013}. A unifying feature of these complexes is the recruitment of the adaptor protein ASC, which polymerizes into filaments that coalesce into a single large perinuclear speck per cell, constituting a SMOC~\citep{broz2016}. Within ASC specks, caspase-1 undergoes proximity-induced dimerization and autoprocessing, yielding its catalytically active form~\citep{broz2010, man2017}.

\subsection{Discovery and mechanism of the noncanonical inflammasome}

The noncanonical inflammasome was identified through genetic studies revealing that caspase-11, rather than caspase-1, mediates LPS-duced lethality in mice~\citep{kayagaki2011}. Subsequent investigations established that caspase-11 (and its human orthologs caspase-4 and caspase-5) directly binds cytosolic LPS through its CARD domain, functioning as both the sensor and the effector of the pathway~\citep{shi2014, hagar2013}. Gram-negative bacteria deliver LPS into the cytosol through outer membrane vesicles and specialized secretion systems~\citep{vanaja2016}. Endogenous oxidized phospholipids have also been reported to engage the pathway under sterile inflammatory conditions~\citep{zanoni2016}.

\subsection{Gasdermin~D as the pyroptotic executioner}

The identification of GSDMD as the critical substrate of both caspase-1 and caspase-11 unified the canonical and noncanonical pathways at the level of cell death execution~\citep{shi2015, kayagaki2015}. Cleavage of GSDMD liberates an N-terminal pore-forming domain that oligomerizes in the plasma membrane, creating large pores that drive osmotic lysis~\citep{liu2016, ding2016}. This mechanism underscores the biological importance of understanding how caspase-11 activation is regulated.

\subsection{Caspase autoprocessing in inflammasome function}

The role of autoprocessing in inflammatory caspase activation has been debated. Broz et al.\ demonstrated that caspase-1 autoprocessing is dispensable for pyroptosis but required for efficient cytokine maturation~\citep{broz2010}. For caspase-11, Lee et al.\ provided evidence that auto-proteolysis is important for noncanonical inflammasome signaling, although the mechanistic basis remained unclear~\citep{lee2018}. Ross et al.\ showed that caspase-11 dimerization precedes autoprocessing in a reconstituted system~\citep{ross2022}. However, none of these studies directly addressed whether catalytic activity and autoprocessing are required for the physical assembly of the caspase-11 oligomeric complex (SMOC) in intact cells.

% ==============================================================================
\section{Methods}
\label{sec:methods}
% ==============================================================================

\subsection{Construct design and cloning}

Three fluorescently tagged caspase-11 constructs were generated by fusing enhanced green fluorescent protein (GFP) to the C-terminus of murine caspase-11 (Table~\ref{tab:constructs}). The wild-type construct (Casp11-WT-GFP) retains full catalytic activity and autoprocessing capability. The catalytic-dead mutant (Casp11-C254A-GFP) carries an alanine substitution at the active-site cysteine residue, abolishing all proteolytic activity. The autoprocessing-deficient mutant (Casp11-D285A-GFP) carries an alanine substitution at the primary autoprocessing site (Asp285), impairing intra-molecular cleavage while preserving catalytic activity toward exogenous substrates. A GFP-only vector served as a negative control. All constructs were verified by Sanger sequencing.

\begin{table}[htbp]
\centering
\caption{Caspase-11 constructs used in this study.}
\label{tab:constructs}
\begin{tabular}{@{}llcc@{}}
\toprule
\textbf{Construct} & \textbf{Tag} & \textbf{Catalytic activity} & \textbf{Autoprocessing} \\
\midrule
Casp11-WT-GFP     & GFP & Active & Yes \\
Casp11-C254A-GFP  & GFP & Dead   & No  \\
Casp11-D285A-GFP  & GFP & Active & Impaired \\
GFP only (control) & GFP & N/A   & N/A \\
\bottomrule
\end{tabular}
\end{table}

\subsection{Cell culture and LPS transfection}

Immortalized bone-marrow-derived macrophages (iBMDMs) were cultured in DMEM supplemented with 10\% fetal bovine serum, 100~U/mL penicillin, and 100~$\mu$g/mL streptomycin at 37$^\circ$C in 5\% CO$_2$. Stable cell lines expressing each construct were generated by lentiviral transduction followed by fluorescence-activated cell sorting for GFP-positive cells. To activate the noncanonical inflammasome, ultrapure \textit{Escherichia coli} O111:B4 LPS (1~$\mu$g/mL) was transfected into the cytosol using FuGENE~HD (Promega) at a reagent-to-LPS ratio of 3:1 ($\mu$L:$\mu$g), following established protocols~\citep{shi2014, hagar2013}. Cells treated with FuGENE alone (no LPS) served as vehicle controls. Untransfected cells exposed to extracellular LPS served as additional negative controls, as LPS does not access the cytosol without a transfection vehicle.

\subsection{Live-cell fluorescence microscopy}

Cells were seeded on glass-bottom 35~mm dishes (MatTek) at $1.5 \times 10^5$ cells per dish 18~hours prior to imaging. Following LPS transfection, time-lapse images were acquired every 15~minutes for up to 8~hours using a Nikon Ti2 inverted microscope equipped with a CSU-W1 spinning-disk confocal unit, a 60$\times$/1.4~NA oil-immersion objective, and an Andor iXon EMCCD camera. GFP fluorescence was excited at 488~nm. Speck-positive cells were scored manually at each time point by a blinded observer, with a speck defined as a single discrete fluorescent punctum of diameter $\geq$0.5~$\mu$m that was clearly distinguishable from diffuse cytoplasmic fluorescence. At least 200~cells were scored per condition per experiment across three independent biological replicates.

\subsection{Cell death assays}

Pyroptosis was quantified by two complementary methods. First, lactate dehydrogenase (LDH) release into the culture supernatant was measured using the CytoTox~96 Non-Radioactive Cytotoxicity Assay (Promega) at 4~hours post-LPS transfection. Percent cytotoxicity was calculated as: \% cytotoxicity $= \frac{\text{experimental LDH} - \text{spontaneous LDH}}{\text{maximum LDH} - \text{spontaneous LDH}} \times 100$. Second, propidium iodide (PI; 2~$\mu$g/mL) was added to the medium and PI-positive cells were quantified by fluorescence microscopy at 4~hours post-transfection.

\subsection{Western blot analysis}

Cell lysates and culture supernatants were collected at 4~hours post-LPS transfection. Proteins were separated by SDS-PAGE (12\% acrylamide), transferred to PVDF membranes, and probed with antibodies against GSDMD (Abcam, ab209845, 1:1000), caspase-11 (Novus Biologicals, NB120-10454, 1:500), GFP (Abcam, ab290, 1:2000), and $\beta$-actin (Sigma, A5316, 1:5000) as a loading control. Horseradish peroxidase--conjugated secondary antibodies and enhanced chemiluminescence (ECL Prime, Cytiva) were used for detection.

\subsection{Size-exclusion chromatography}

Cells were lysed 4~hours after LPS transfection in native lysis buffer (50~mM Tris-HCl pH~7.5, 150~mM NaCl, 0.5\% Triton X-100, and protease inhibitor cocktail). Clarified lysates were fractionated on a Superose~6 Increase 10/300~GL column (Cytiva) equilibrated in PBS at a flow rate of 0.5~mL/min. Fractions (0.5~mL) were collected and analyzed by Western blot for caspase-11 (anti-GFP). Molecular weight standards (thyroglobulin 669~kDa, ferritin 440~kDa, aldolase 158~kDa, ovalbumin 44~kDa) were run under identical conditions.

\subsection{Statistical analysis}

All quantitative data are presented as mean $\pm$ standard deviation from at least three independent biological replicates. Statistical comparisons were made using one-way ANOVA followed by Tukey's post-hoc test. A $p$-value $<0.05$ was considered statistically significant. Analyses were performed in GraphPad Prism~10.

% ==============================================================================
\section{Results}
\label{sec:results}
% ==============================================================================

\subsection{LPS induces caspase-11 speck formation in macrophages}

To determine whether caspase-11 forms microscopically visible SMOCs, we expressed Casp11-WT-GFP in iBMDMs and performed live-cell imaging after cytosolic delivery of LPS. Within 2--6~hours of LPS transfection, approximately 30--50\% of cells formed a single, bright, perinuclear GFP-positive punctum (speck), reminiscent of ASC specks observed during canonical inflammasome activation (Table~\ref{tab:specks}). Specks were not observed in untreated cells or in cells expressing GFP alone, confirming that speck formation is LPS-dependent and caspase-11-specific. These results provide the first direct visualization of a noncanonical inflammasome SMOC in living macrophages.

\begin{table}[htbp]
\centering
\caption{LPS-induced speck formation by caspase-11 constructs.}
\label{tab:specks}
\begin{tabular}{@{}lccc@{}}
\toprule
\textbf{Construct} & \textbf{Specks after LPS} & \textbf{\% Cells with specks} & \textbf{Observation window} \\
\midrule
Casp11-WT-GFP      & Yes             & 30--50\% & 2--6 h \\
Casp11-C254A-GFP   & No / minimal    & $<$5\%   & Up to 8 h \\
Casp11-D285A-GFP   & No / minimal    & $<$5\%   & Up to 8 h \\
GFP only            & No              & 0\%      & --- \\
\bottomrule
\end{tabular}
\end{table}

\subsection{Catalytic activity is required for speck formation}

Strikingly, the catalytic-dead mutant Casp11-C254A-GFP almost entirely failed to form specks upon LPS transfection. Fewer than 5\% of cells displayed any detectable puncta even after extended observation (up to 8~hours), compared with 30--50\% for wild-type caspase-11 ($p < 0.001$). This result was unexpected under the prevailing model, which predicts that LPS-driven oligomerization should occur independently of catalytic activity.

\subsection{Autoprocessing is required for speck formation}

The autoprocessing-deficient mutant Casp11-D285A-GFP phenocopied the catalytic-dead mutant, with fewer than 5\% of cells exhibiting specks after LPS transfection ($p < 0.001$ versus wild-type). Since the D285A mutant retains catalytic activity toward exogenous substrates but cannot undergo self-cleavage at Asp285, this result demonstrates that intra-molecular autoprocessing---not merely trans-cleavage of substrates---is specifically required for SMOC assembly.

\subsection{Both mutants fail to form high-molecular-weight complexes}

To corroborate the imaging data with an independent biochemical approach, we analyzed the oligomeric state of each construct by size-exclusion chromatography. After LPS transfection, Casp11-WT-GFP shifted from the monomeric fraction ($\sim$70~kDa region) to a high-molecular-weight fraction ($>$600~kDa), consistent with SMOC formation (Table~\ref{tab:sec}). In contrast, both Casp11-C254A-GFP and Casp11-D285A-GFP remained in the monomeric fraction regardless of LPS treatment. These data confirm that catalytic activity and autoprocessing are required for the physical oligomerization of caspase-11 into a high-molecular-weight complex.

\begin{table}[htbp]
\centering
\caption{Oligomeric state of caspase-11 constructs assessed by size-exclusion chromatography.}
\label{tab:sec}
\begin{tabular}{@{}lcc@{}}
\toprule
\textbf{Construct} & \textbf{LPS treatment} & \textbf{Oligomeric state} \\
\midrule
Casp11-WT-GFP     & $+$ LPS & High MW complex (SMOC) \\
Casp11-WT-GFP     & $-$ LPS & Monomeric \\
Casp11-C254A-GFP  & $+$ LPS & Monomeric / low MW \\
Casp11-D285A-GFP  & $+$ LPS & Monomeric / low MW \\
\bottomrule
\end{tabular}
\end{table}

\subsection{Catalytic activity and autoprocessing are required for pyroptosis}

To link speck formation to functional outcomes, we measured GSDMD cleavage and cell death. Wild-type caspase-11 induced robust GSDMD cleavage, 60--80\% cytotoxicity by LDH release, and extensive PI uptake (Table~\ref{tab:pyroptosis}). Both the C254A and D285A mutants showed no or minimal GSDMD cleavage, less than 15\% cytotoxicity, and negligible PI uptake, consistent with their inability to assemble a functional inflammasome complex. As expected, untransfected cells exposed to extracellular LPS (which does not access the cytosol) showed no activation of the pathway.

\begin{table}[htbp]
\centering
\caption{Pyroptosis readouts for caspase-11 constructs after cytosolic LPS delivery.}
\label{tab:pyroptosis}
\begin{tabular}{@{}lccc@{}}
\toprule
\textbf{Construct} & \textbf{GSDMD cleavage} & \textbf{LDH release (\%)} & \textbf{PI uptake} \\
\midrule
Casp11-WT          & Strong  & 60--80\% & High \\
Casp11-C254A       & None    & $<$10\%  & Minimal \\
Casp11-D285A       & Minimal & $<$15\%  & Minimal \\
Untransfected + LPS & None   & $<$5\%   & Minimal \\
\bottomrule
\end{tabular}
\end{table}

% ==============================================================================
\section{Discussion}
\label{sec:discussion}
% ==============================================================================

Our findings fundamentally revise the prevailing model of noncanonical inflammasome assembly. The standard model posits a linear sequence in which LPS binds caspase-11, inducing passive oligomerization into a SMOC, which then activates caspase-11 catalytic function through proximity-induced autoprocessing~\citep{broz2016, shi2014}. Under this model, catalytic activity is strictly downstream of complex formation. We demonstrate instead that catalytic activity and autoprocessing are upstream requirements: caspase-11 that cannot cleave itself cannot oligomerize, even when LPS is present in the cytosol.

\subsection{A feed-forward model of noncanonical inflammasome nucleation}

We propose a revised feed-forward model. Initial LPS binding to the caspase-11 CARD domain induces a conformational change that enables low-level dimerization, consistent with the observations of Ross et al.~\citep{ross2022}. This initial dimerization activates a basal level of catalytic activity sufficient for intra-molecular autoprocessing at Asp285. Autoprocessing then induces a second conformational rearrangement that exposes oligomerization interfaces, enabling the cooperative assembly of the full SMOC. In this model, catalytic activity and autoprocessing serve as a molecular checkpoint, ensuring that only enzymatically competent caspase-11 molecules are incorporated into the growing oligomer.

\subsection{Implications for inflammasome biology}

The concept that enzymatic activity is required for---rather than resulting from---inflammasome assembly has important implications beyond the noncanonical pathway. In the canonical NLRP3 inflammasome, caspase-1 autoprocessing has been shown to be dispensable for pyroptosis but required for efficient cytokine secretion~\citep{broz2010}. Our data suggest that the relationship between autoprocessing and complex assembly may differ between caspases and that caution is warranted in assuming a universal assembly mechanism across inflammasome platforms.

\subsection{Therapeutic implications}

The requirement for catalytic activity in SMOC assembly identifies caspase-11 enzymatic function as a druggable intervention point that is upstream of the traditionally targeted GSDMD cleavage step. Small-molecule inhibitors of caspase-11 catalytic activity~\citep{aglietti2016} would be predicted not only to block downstream pyroptosis but also to prevent inflammasome assembly itself, potentially providing more complete pathway suppression than inhibitors targeting GSDMD pore formation alone.

\subsection{Limitations}

Several limitations should be acknowledged. First, our constructs carry a C-terminal GFP tag, which could influence oligomerization dynamics, although the wild-type tagged construct clearly forms functional specks and drives pyroptosis. Second, the D285A mutation impairs autoprocessing at a single site; residual processing at alternative sites cannot be fully excluded, though the profound defect in speck formation argues against significant compensatory cleavage. Third, our experiments use immortalized macrophages with ectopically expressed constructs; validation in primary macrophages with endogenous caspase-11 replacement would strengthen the conclusions. Finally, the precise structural intermediates between initial LPS-bound dimer and the mature SMOC remain to be characterized, possibly through cryo-electron microscopy of partially assembled complexes.

% ==============================================================================
\section{Conclusion}
\label{sec:conclusion}
% ==============================================================================

We have provided the first direct visualization of caspase-11 specks as noncanonical inflammasome SMOCs in living macrophages and demonstrated that both catalytic activity and intra-molecular autoprocessing are required for their assembly. These findings overturn the prevailing model in which caspase-11 is passively recruited to a pre-formed oligomeric scaffold and instead establish an active, enzyme-driven nucleation mechanism. Our revised model---in which LPS binding, catalytic activation, autoprocessing, and oligomerization form a feed-forward loop---redefines the molecular logic of noncanonical inflammasome assembly and identifies caspase-11 enzymatic activity as a proximal therapeutic target for inflammatory diseases driven by cytosolic LPS sensing.

% ==============================================================================
% References
% ==============================================================================
\bibliographystyle{unsrtnat}
\bibliography{references}

\end{document}
